\chapter{The Gaussian Chain: Exact Score, Belief Propagation, and
Locality}\label{ch:gaussian}

This chapter answers RQ1, RQ2, and the Gaussian half of RQ4. For the
stationary AR(1) chain under the OU channel we compute the joint score
exactly by two independent routes (linear algebra and Bayesian smoothing),
characterise how its coupling structure evolves with diffusion time,
derive the belief-propagation computation of the same score and identify
it with the Kalman filter and smoother, and quantify how rapidly the
influence of remote observations decays. Everything in this chapter is
exact unless explicitly marked as a numerical check; each boxed formula is
covered by a named check of the audit suite (Appendix~\ref{app:repro}).

\section{The joint law and the object of study}\label{sec:g-model}

A data point is a trajectory $a = (a_0, \dots, a_{K-1})$ of the stationary
AR(1) chain \eqref{eq:ar1} with $\sigma_\eta^2 = 1 - \alpha^2$. Each frame
passes independently through the OU channel \eqref{eq:ou-channel} for a
time $t$,
\begin{equation}
  x_k \;=\; e^{-t} a_k + \sqrt{\Deltat}\;\xi_k,
  \qquad \xi_k \sim \Nd(0,1)\ \text{i.i.d.},
  \label{eq:g-channel}
\end{equation}
and the object of study is the joint score of the noisy sequence,
\begin{equation}
  \score(x, t) \;=\; \nabla_x \log P_t(x) \;\in\; \R^K,
  \qquad
  P_t(x) = \int_{\R^K} P_0(a) \prod_{k=0}^{K-1}
    \Nd\big(x_k;\, e^{-t}a_k,\, \Deltat\big)\, \dd a .
  \label{eq:g-object}
\end{equation}
The gradient is taken with respect to all $K$ noisy frames simultaneously,
under the law of the whole noisy trajectory
(Section~\ref{sec:tweedie}).

\section{The clean law: dense covariance, sparse precision}
\label{sec:g-clean}

The chain factorisation \eqref{eq:chain-factorisation} gives the clean
joint density
\begin{equation}
  P_0(a)
  = \Nd(a_0; 0, 1) \prod_{k=0}^{K-2}
    \Nd\big(a_{k+1};\, \alpha a_k,\, \sigma_\eta^2\big)
  = \Nd(a;\, 0,\, \Sigz).
  \label{eq:g-P0}
\end{equation}

\begin{proposition}[Stationary covariance]\label{prop:g-Sigma0}
$(\Sigz)_{ij} = \alpha^{|i-j|}$: a dense, symmetric Toeplitz matrix with
geometrically decaying off-diagonals, $\gamma(d) = \alpha^{|d|}$.
\end{proposition}

\begin{proposition}[Tridiagonal precision]\label{prop:g-Q0}
$\Qz := \Sigz^{-1}$ is exactly tridiagonal:
\begin{equation}
  (\Qz)_{kk} =
  \begin{cases}
    \dfrac{1}{\sigma_\eta^2}, & k \in \{0, K-1\},\\[8pt]
    \dfrac{1+\alpha^2}{\sigma_\eta^2}, & 0 < k < K-1,
  \end{cases}
  \qquad
  (\Qz)_{k, k\pm 1} = -\frac{\alpha}{\sigma_\eta^2},
  \label{eq:g-Q0}
\end{equation}
and every entry at distance $\geq 2$ vanishes identically.
\end{proposition}

\begin{toolbox}{Deriving the tridiagonal precision from the log-density}
For a centred Gaussian, the precision is the Hessian of the negative
log-density. From \eqref{eq:g-P0},
\begin{equation*}
  -2 \log P_0(a)
  = \frac{a_0^2}{\sigma_0^2}
  + \sum_{k=0}^{K-2} \frac{(a_{k+1} - \alpha a_k)^2}{\sigma_\eta^2}
  + \text{const},
  \qquad \sigma_0^2 = 1 .
\end{equation*}
Read off the quadratic form term by term. The only term coupling $a_k$ and
$a_{k+1}$ is the transition $(a_{k+1} - \alpha a_k)^2/\sigma_\eta^2$, with
cross coefficient $-2\alpha/\sigma_\eta^2$, hence precision entry
$-\alpha/\sigma_\eta^2$; no term couples variables at distance $\geq 2$,
so those entries are exactly zero. An interior $a_k$ appears in its left
transition with coefficient $1/\sigma_\eta^2$ and in its right transition
with $\alpha^2/\sigma_\eta^2$, summing to $(1+\alpha^2)/\sigma_\eta^2$. At
the left boundary the prior term $1/\sigma_0^2 = (1-\alpha^2)/\sigma_\eta^2$
plus the right transition $\alpha^2/\sigma_\eta^2$ give $1/\sigma_\eta^2$;
the right boundary is symmetric. \hfill$\square$
\end{toolbox}

\begin{remark}[Tridiagonality is Markovianity, not Gaussianity]
\label{rmk:g-markov-sparsity}
A zero in the precision matrix is a conditional independence given all
other variables; by the Hammersley--Clifford theorem
(Section~\ref{sec:mrf}), factorisation over the cliques of a graph and
conditional independence with respect to that graph are the same property.
Gaussianity entered the derivation only to identify the precision with the
Hessian of $-\log P_0$; the support of that Hessian, nearest neighbours
only, comes from the chain factorisation alone. Any Markov chain with
smooth transitions has this sparsity pattern; the Gaussian case makes the
entries constants.
\end{remark}

The covariance is dense (every frame correlates with every other) while
the precision is tridiagonal: correlation is global, conditional
dependence is local. The score is governed by the second object.

\section{The score, two ways}\label{sec:g-score}

\subsection{Linear algebra}

\begin{proposition}[Noised law]\label{prop:g-Sigt}
$P_t(x) = \Nd(x;\, 0,\, \Sigt)$ with
\begin{equation}
  \boxed{\;\Sigt \;=\; e^{-2t}\, \Sigz \;+\; \Deltat\, I_K\;}
  \label{eq:g-Sigt}
\end{equation}
\end{proposition}
This is \eqref{eq:ou-cov-map}. The diagonal is
$e^{-2t} + \Deltat = 1$ at every $t$ (variance preservation); only the
off-diagonal entries, $\alpha^{|i-j|} e^{-2t}$, depend on the diffusion
time, because the noise of different frames is independent. Since
$\log P_t(x) = -\tfrac{1}{2} x^\T \Sigt^{-1} x + \text{const}$,
\begin{equation}
  \boxed{\;\score(x, t) \;=\; -\,\Qt\, x,
  \qquad \Qt := \Sigt^{-1}.\;}
  \label{eq:g-score-matrix}
\end{equation}
The joint score is linear in $x$, the hallmark of Gaussianity, and its
entire structure is the noisy precision matrix $\Qt$: which frames couple,
how strongly, and at which $t$, are questions about this one matrix,
answered in Sections~\ref{sec:g-lifecycle} and~\ref{sec:g-bulk}.

\subsection{Bayesian inference}

The score--posterior identity \eqref{eq:tweedie-joint} expresses the same
score through the posterior mean of the clean chain,
\begin{equation}
  \score_k(x, t)
  = \frac{e^{-t}\, \E[a_k \,|\, x] - x_k}{\Deltat} ,
  \label{eq:g-tweedie}
\end{equation}
and the posterior is Gaussian with parameters used throughout:

\begin{proposition}[Posterior of the clean chain]\label{prop:g-posterior}
$P(a \,|\, x) = \Nd(a;\, m,\, J^{-1})$ with
\begin{equation}
  \boxed{\;
  J \;=\; \frac{e^{-2t}}{\Deltat}\, I \;+\; \Qz ,
  \qquad
  J\, m \;=\; h := \frac{e^{-t}}{\Deltat}\, x .
  \;}
  \label{eq:g-posterior}
\end{equation}
\end{proposition}

\begin{toolbox}{Posterior precision and field, by completing the square}
Bayes' rule with the likelihood \eqref{eq:g-channel}:
\begin{align*}
  -2 \log P(a | x)
  &= a^\T \Qz\, a
  + \sum_{k} \frac{(x_k - e^{-t} a_k)^2}{\Deltat}
  + \text{const} \\
  &= a^\T \Big( \Qz + \tfrac{e^{-2t}}{\Deltat} I \Big) a
  - 2\, \tfrac{e^{-t}}{\Deltat}\, x^\T a
  + \text{const}.
\end{align*}
Matching against the generic form $a^\T J a - 2 h^\T a$ gives
\eqref{eq:g-posterior}. Two structural observations. (i) The likelihood is
diagonal in $a$: observing each frame through its own channel adds the
channel precision $e^{-2t}/\Deltat$ to every diagonal entry of the prior
precision and touches nothing else. (ii) Consequently $J$ is tridiagonal
at every $t$: the \emph{posterior} of the clean chain remains a
nearest-neighbour Markov field; what fills up with $t$ is the precision
$\Qt$ of the noisy marginal law. Keeping $J$ (tridiagonal, posterior) and
$\Qt$ (dense, noisy marginal) distinct prevents much confusion.
\hfill$\square$
\end{toolbox}

Substituting $m = J^{-1} h$ into \eqref{eq:g-tweedie} must return
$-\Qt x$; the operator identity behind this is Woodbury's
(Appendix~\ref{app:gauss}), and the audit suite checks the two routes
agree to $2.7 \times 10^{-15}$ across sweeps of $(K, \alpha, t)$. The
agreement is an end-to-end check of two computations built from different
formulas.

\subsection{The smallest example, $K = 2$}\label{sec:g-K2}

For $K = 2$, with $r := \alpha\, e^{-2t}$,
\begin{equation}
  \Sigt = \begin{pmatrix} 1 & r \\ r & 1 \end{pmatrix}
  \qquad\Longrightarrow\qquad
  \score_0(x, t) = -\,\frac{x_0 - r\, x_1}{1 - r^2},
  \qquad
  \score_1(x, t) = -\,\frac{x_1 - r\, x_0}{1 - r^2}.
  \label{eq:g-K2-score}
\end{equation}
By contrast, each $x_k$ is marginally $\Nd(0,1)$ at every $t$ (variance
preservation), so the per-frame marginal score is $-x_k$: independent of
$t$ and of the other frame, it carries no information about the dynamics.
The joint score contains the cross term $-r x_1/(1-r^2)$: the denoiser of
frame $0$ uses evidence from frame $1$ with a weight that equals $\alpha$
at $t = 0$ and decays as $e^{-2t}$. This is the coupling structure whose
full lifecycle the next section derives, in miniature.

\section{The spectral form and the lifecycle of the precision}
\label{sec:g-lifecycle}

How does $\Qt$ interpolate between the tridiagonal $\Qz$ at $t = 0$ and
the identity at $t = \infty$? The analysis rests on the spectral theorem
and its affine-family corollary (Section~\ref{sec:matrix-tools}): affine
images of a symmetric matrix share its eigenbasis,
$(cA + dI)\, u_i = (c\,\omega_i + d)\, u_i$. Applying this to the
forward map \eqref{eq:g-Sigt}, which is affine in $\Sigz$, gives the
representation from which everything in this section follows.

\begin{theorem}[Spectral form]\label{thm:g-three-forms}
Let $\Sigz = U \operatorname{diag}(\omega_i) U^\T$. Then for every
$t \geq 0$,
\begin{equation}
  \boxed{\;
  \Sigt = U \operatorname{diag}\big( \lambda_i(t) \big) U^\T,
  \qquad
  \Qt = U \operatorname{diag}\!\Big( \frac{1}{\lambda_i(t)} \Big) U^\T,
  \qquad
  \lambda_i(t) = e^{-2t} \omega_i + \Deltat .
  \;}
  \label{eq:g-Qt-spec}
\end{equation}
\end{theorem}

Diffusion never rotates the problem: the eigenvectors of $\Sigt$ and
$\Qt$ are those of $\Sigz$, frozen for all $t$, while each eigenvalue
flows along the affine path $\omega_i \mapsto e^{-2t}\omega_i + \Deltat$,
monotonically towards $1$. Because $\Sigz$ is symmetric Toeplitz, its
eigenvectors are asymptotically the Fourier modes of the sequence
(Section~\ref{sec:matrix-tools}): in frequency space the denoising problem is $K$
independent scalar problems at every diffusion time. Two computational
corollaries are recorded for later use: factoring $\Deltat$ out of
\eqref{eq:g-Sigt} gives the SNR form
\begin{equation}
  \Qt = \tfrac{1}{\Deltat} \big( I + \gamma_t\, \Sigz \big)^{-1},
  \qquad \gamma_t := e^{-2t}/\Deltat ,
  \label{eq:g-Qt-snr}
\end{equation}
and multiplying \eqref{eq:g-Sigt} by $\Qz$ and inverting gives the
resolvent form
\begin{equation}
  \Qt = e^{2t}\, \Qz \big( I + c_t \Qz \big)^{-1}
      = e^{2t} \sum_{m \geq 0} (-c_t)^m\, \Qz^{\,m+1},
  \qquad c_t := e^{2t} - 1 ,
  \label{eq:g-Qt-res}
\end{equation}
whose Neumann series converges for $\|c_t \Qz\| < 1$, i.e.\ at small $t$.

\subsection{Losing tridiagonality: when, how, and how much}
\label{sec:g-tridiag-loss}

Sparsity is a property of a matrix in a basis. In the eigenbasis the
problem is diagonal at all times; in the frame basis, the operationally
meaningful one, the zeros of $\Qz$ are a cancellation that diffusion
destroys. The destruction is quantifiable, and it answers three questions:
when tridiagonality is lost, how, and how much structure survives at
each $t$.

\paragraph{When: immediately.} By \eqref{eq:g-Qt-spec} every entry of
$\Qt$ is analytic in $t$. The distance-$d$ entries, identically zero at
$t = 0$ for $d \geq 2$, are non-zero for every $t > 0$: no exact zero
survives any positive diffusion time.

\paragraph{How: one diagonal per order in $t$.} The loss is nevertheless
graded, and the resolvent form gives the exact rate at which each diagonal
switches on.

\begin{theorem}[Band-fill leading order]\label{thm:g-bandfill}
For small $t$ and frame distance $d \geq 1$,
\begin{equation}
  \boxed{\;
  (\Qt)_{i, i+d}
  \;=\; (-1)^{d-1}\, (2t)^{d-1}\, \big( \Qz^{\,d} \big)_{i, i+d}
  \;+\; O(t^{d}) .
  \;}
  \label{eq:g-bandfill}
\end{equation}
\end{theorem}

\begin{toolbox}{Band fill from the Neumann series}
In the resolvent form \eqref{eq:g-Qt-res}, substitute
$c_t = 2t + 2t^2 + O(t^3)$ and $e^{2t} = 1 + O(t)$:
\begin{equation*}
  \Qt = \big(1 + O(t)\big)
  \Big[ \Qz - c_t \Qz^2 + c_t^2 \Qz^3 - \cdots \Big].
\end{equation*}
$\Qz$ is tridiagonal, so $\Qz^{\,m}$ has bandwidth $m$: the entry at
distance $d$ receives its first contribution from the single term $m = d$,
with prefactor $(-c_t)^{d-1} = (-1)^{d-1}(2t)^{d-1}(1 + O(t))$; every
later term is $O(t^d)$. Each additional frame of coupling range costs one
power of $t$: the band fills outward, one diagonal per order.
\hfill$\square$
\end{toolbox}

\begin{remark}[The prefactor carries no factorial, and why that is checked]
\label{rmk:g-correction}
Equation~\eqref{eq:g-bandfill} has no $1/(d-1)!$: the leading distance-$d$
term is a single Neumann term, with no combinatorial sum behind it. The
distinction is easy to miss and worth verifying numerically, because slope
fits in $\log t$ are blind to constant prefactors --- a wrong constant
survives every check on the exponent. The audit therefore tests the
coefficient itself, and it separates the two candidates cleanly: the
no-factorial form matches the measured $(\Qt)_{i,i+d}/t^{d-1}$ to better
than $0.9\%$ for every $d \in \{1,\dots,5\}$ (at $\alpha = 0.9$, $K = 25$),
while a factorial version would be wrong by $86\%$ already at $d = 3$. Every
prefactor in this thesis is audited on the same grounds, not only every
exponent.
\end{remark}

\begin{figure}[t]\centering
\includegraphics[width=0.94\linewidth]{fig_band_fill.pdf}
\caption[The band-fill law for the posterior precision]{Band-fill law. Measured $|(\Qt)_{i,i+d}|$ against $t$ (dots) and
the leading order \eqref{eq:g-bandfill} (solid lines, no fitted
parameters), for $\alpha = 0.9$, $K = 25$, $d = 1, \dots, 5$. The slopes
are $d - 1$; agreement at small $t$ is at the percent level for every
$d$.}
\label{fig:g-bandfill}
\end{figure}

\paragraph{How much: the off-tridiagonal mass.}
Define the off-tridiagonal Frobenius mass of $\Qt$: the norm of everything
outside the band that Markov structure would allow. At small $t$ the
band-fill law makes it rise from zero, dominated by the $d = 2$ diagonal,
linearly in $t$. At large $t$, expanding the inverse of
$\Sigt = I + e^{-2t}(\Sigz - I)$ gives
\begin{equation}
  \Qt = I - e^{-2t}(\Sigz - I) + e^{-4t}(\Sigz - I)^2 + O(e^{-6t}),
  \label{eq:g-larget}
\end{equation}
so $\| \Qt - I \|_F \simeq e^{-2t} \|\Sigz - I\|_F$: all structure,
off-tridiagonal and tridiagonal alike, decays at the rate $e^{-2t}$, and
the score approaches the isotropic $-x$. In between, the off-tridiagonal
mass peaks (Figure~\ref{fig:g-tridiagloss}): the noisy sequence is most
non-Markov at intermediate diffusion times. The lifecycle in one line:
tridiagonal at $t = 0$; densifying at rate $t^{d-1}$ per diagonal $d$;
maximally coupled at intermediate $t$; isotropic as $t \to \infty$ at rate
$e^{-2t}$.

\begin{figure}[t]\centering
\includegraphics[width=0.94\linewidth]{fig_tridiag_loss.pdf}
\caption[Loss and recovery of tridiagonality with diffusion time]{Structure across diffusion time ($\alpha = 0.8$, $K = 40$).
$\|\Sigt - \Sigz\|_F$ (navy) grows monotonically as the covariance is
driven to the identity; the off-tridiagonal Frobenius mass of $\Qt$ (red)
rises from zero, peaks at intermediate $t$, and decays to zero: the band
fills, then empties.}
\label{fig:g-tridiagloss}
\end{figure}

\section{Belief propagation on the chain}\label{sec:bp-chain}\label{sec:bp-problem}

The matrix route inverts a $K \times K$ matrix, a global computation
over the whole sequence. The score--posterior identity says the same object is a posterior
expectation, and the posterior has a structure that supports local
computation. This section develops that computation once, in full.

\subsection{The factor graph of the posterior}

For any Markov-chain prior and per-frame OU evidence, Bayes' rule gives
\begin{equation}
  p(a \,|\, x) \;\propto\;
  \psi_{\mathrm{pr}}(a_0)
  \prod_{k=0}^{K-2} \psi^{(k)}_{\mathrm{tr}}(a_k, a_{k+1})
  \prod_{k=0}^{K-1} \psi^{(k)}_{\mathrm{ob}}(a_k),
  \qquad
  \psi^{(k)}_{\mathrm{ob}}(a_k) = \Nd\big(x_k;\, e^{-t}a_k,\, \Deltat\big),
  \label{eq:bp-posterior}
\end{equation}
a hidden Markov model \citep{rabiner1989tutorial}: a latent chain observed
through independent per-frame noise. By \eqref{eq:tweedie-joint}, the
$k$-th score component requires exactly the smoothing marginal
$\E[a_k | x]$ of \eqref{eq:bp-posterior}.

The factor graph of this posterior (Section~\ref{sec:factor-graphs})
has $K$ variable nodes, $2K$
factor nodes (one prior, $K{-}1$ transitions, $K$ observations; the data
$x_k$ are constants clamped inside their factors) and $3K - 1$ edges. A
connected graph with one fewer edge than nodes is a tree: here, a spine of
variables and transition factors with one evidence leaf per variable.

\begin{proposition}\label{prop:chain-tree}
For any Markov-chain prior and any per-frame observation model, the
posterior $p(a \,|\, x)$ factorises over a tree-shaped factor graph, and
all $K$ smoothing marginals are computed exactly by one forward and one
backward sweep of sum--product messages.
\end{proposition}

Proposition~\ref{prop:chain-tree} counts sweeps, not work: it says
nothing about how the messages are represented, and for continuous
variables each message is a function
(Section~\ref{sec:sum-product}). Keeping tree-exactness of the
algorithm (a graph property) separate from representability of the
messages (a property of the model's function class) is the conceptual
hinge of this thesis; the Gaussian case closes the gap, and
Chapter~\ref{ch:laplace} measures what happens when it no longer does.

\subsection{The sum--product recursion, with the bookkeeping explicit}

Specialising the sum--product rules \eqref{eq:bp-rules} to this tree,
the schedule collapses into one forward and one backward sweep
(Theorem~\ref{thm:tree-exact}), and the bookkeeping deserves care.

\begin{proposition}[Chain sweeps, Convention A]\label{prop:bp-convA}
Define messages
\begin{align}
  \mu_{\to 0}(a_0) &:= \psi_{\mathrm{pr}}(a_0),
  &
  \mu_{\to k+1}(a_{k+1}) &:= \int
    \psi^{(k)}_{\mathrm{tr}}(a_k, a_{k+1})\,
    \psi^{(k)}_{\mathrm{ob}}(a_k)\,
    \mu_{\to k}(a_k)\; \dd a_k ,
  \label{eq:bp-fwd}\\
  \mu_{\leftarrow K-1}(a_{K-1}) &:= 1,
  &
  \mu_{\leftarrow k-1}(a_{k-1}) &:= \int
    \psi^{(k-1)}_{\mathrm{tr}}(a_{k-1}, a_k)\,
    \psi^{(k)}_{\mathrm{ob}}(a_k)\,
    \mu_{\leftarrow k}(a_k)\; \dd a_k .
  \label{eq:bp-bwd}
\end{align}
Then for every $k$,
\begin{equation}
  p(a_k \,|\, x) \;\propto\;
  \mu_{\to k}(a_k)\;
  \psi^{(k)}_{\mathrm{ob}}(a_k)\;
  \mu_{\leftarrow k}(a_k) :
  \label{eq:bp-combine}
\end{equation}
the forward message into $a_k$ carries the evidence
$x_0, \dots, x_{k-1}$ only, the backward message carries
$x_{k+1}, \dots, x_{K-1}$ only, and the local evidence enters once, at
combination time.
\end{proposition}

The proof is direct marginalisation: integrating \eqref{eq:bp-posterior}
over all $a_j$, $j \neq k$, splits at $a_k$ by the Markov property into
independent left and right halves, and peeling the innermost variable of
each half yields the recursions.

\begin{remark}[The double-counting trap]\label{rmk:bp-trap}
An alternative convention absorbs the local evidence into the forward
message, $\mu^{\mathrm{B}}_{\to k} = \mu_{\to k} \cdot
\psi^{(k)}_{\mathrm{ob}}$; the product \eqref{eq:bp-combine} then counts
$x_k$ twice and must be corrected by dividing one factor back out. The cost
of getting this wrong is not subtle, which is why it is worth stating: under
Convention A the BP score matches the matrix score to $10^{-14}$, and under the
naive alternative the disagreement is of order one. Every message in this thesis
is therefore stated with its evidence content explicit.
\end{remark}

\section{The Gaussian closure and the Kalman smoother}
\label{sec:bp-gaussian}\label{sec:bp-kalman}

For the Gaussian chain, the two closure lemmas of
Toolbox~\ref{tb:closure} are the only facts needed: the recursions
\eqref{eq:bp-fwd}--\eqref{eq:bp-bwd} perform exactly one product and one
affine marginalisation per step. The local evidence, read as a function
of $a_k$, is itself Gaussian,
$\psi^{(k)}_{\mathrm{ob}}(a_k) \propto
\Nd(a_k;\, e^{t} x_k,\, \Deltat e^{2t})$: a pseudo-observation at the
deconvolved location $e^t x_k$ with precision $e^{-2t}/\Deltat$, the
pinning strength that appears throughout this chapter.

\begin{theorem}[Gaussian closure of BP on the chain]\label{thm:bp-closure}
For the Gaussian AR(1) chain under the OU channel, every BP message and
every belief is exactly Gaussian, at every node, for every
$(K, \alpha, t)$. Each message is carried by two numbers, each update
is a few arithmetic operations, and the two sweeps compute the exact
posterior marginals, hence the exact joint score. Here, and only here,
the sweep count of Proposition~\ref{prop:chain-tree} translates into an
actual computation of fixed, small cost per frame.
\end{theorem}

\begin{proof}[Derivation]
Induction along each sweep. The initial messages are Gaussian (the
prior; the flat message is the zero-precision limit of the family). Each
step of \eqref{eq:bp-fwd}--\eqref{eq:bp-bwd} is one application of
Lemma~1 of Toolbox~\ref{tb:closure} (absorb the local evidence) followed
by one of Lemma~2 (push through the transition); the combination
\eqref{eq:bp-combine} is Lemma~1 twice. Exactness follows from tree
exactness of sum--product (Theorem~\ref{thm:tree-exact}).
\end{proof}

\begin{remark}[Closure is algebra, not a central limit theorem]
\label{rmk:bp-no-clt}
The closure just proved holds at node degree two exactly as it would at
any degree: it is algebra, not an aggregation effect, as discussed in
Section~\ref{sec:gaussian-bp}. The central-limit argument matters only
for the per-node closures of statistical physics
(Section~\ref{sec:amp-intro}), and on the chain those fail in a way that
can be made exact, with a breakdown time $t_c(\alpha)$ and a critical
coupling $\alpha_c = \sqrt{2} - 1$ (Appendix~\ref{app:amp}).
\end{remark}

Executing the Gaussian sweeps in mean--variance arithmetic reproduces,
update by update, the classical algorithms of linear estimation.

\begin{toolbox}{The Kalman filter, derived from the BP sweep}
Write the forward message in Convention A as
$\mu_{\to k} = \Nd(a_k;\, m_{k|k-1},\, P_{k|k-1})$, the predicted state
given past observations. One BP step does two things.

\emph{Update (absorb the evidence; Lemma 1).} Multiplying by the
pseudo-observation $\Nd(a_k; e^t x_k, \Deltat e^{2t})$ gives the filtered
Gaussian $\Nd(a_k; m_{k|k}, P_{k|k})$ with
\begin{align*}
  \frac{1}{P_{k|k}} &= \frac{1}{P_{k|k-1}} + \frac{e^{-2t}}{\Deltat},\\
  m_{k|k} &= P_{k|k} \left( \frac{m_{k|k-1}}{P_{k|k-1}}
        + \frac{e^{-t} x_k}{\Deltat} \right)
  = m_{k|k-1} + G_k \big( e^{t} x_k - m_{k|k-1} \big),
\end{align*}
where $G_k = P_{k|k}\, e^{-2t}/\Deltat \in (0,1)$ is the Kalman gain: the
filtered mean is the prediction corrected by a gain-weighted innovation
\citep{kalman1960new}.

\emph{Predict (push through the dynamics; Lemma 2).}
\begin{equation*}
  m_{k+1|k} = \alpha\, m_{k|k},
  \qquad
  P_{k+1|k} = \alpha^2 P_{k|k} + \sigma_\eta^2 .
\end{equation*}
The forward BP sweep is the Kalman filter, message by message.

\emph{Smoothing (combine both directions).} The belief
\eqref{eq:bp-combine} multiplies the filtered forward Gaussian by the
backward message (Lemma~1 once more), producing the posterior marginal
given all observations, $\Nd(a_k; m_{k|K-1}, P_{k|K-1})$; carrying out
the same algebra for the backward recursion yields the fixed-interval
smoother of \citet{rauch1965maximum}. The two-sweep BP schedule on the
chain is the Kalman filter followed by the RTS smoother, exactly as the
generic correspondence of Section~\ref{sec:gaussian-bp} predicts.
\hfill$\square$
\end{toolbox}

Combining with the score--posterior identity answers RQ2: the $k$-th
component of the joint diffusion score is an affine function of the
Kalman-smoothed estimate of frame $k$,
\begin{equation}
  S_k(x, t)
  \;=\; \frac{e^{-t}\, m_{k|K-1} - x_k}{\Deltat} .
  \label{eq:bp-score-kalman}
\end{equation}
For Markov sequences, score computation is a smoothing problem, and the
classical theory of linear smoothing applies directly.

\paragraph{Numerical agreement of the two routes.} The audit suite runs BP
against the matrix route over $80$ configurations
($K \in \{2,3,5,8,16\}$, $\alpha \in \{0.2, 0.5, 0.9, -0.4\}$,
$t \in \{0.05, 0.3, 1, 3\}$, random $x$): maximum disagreement
$1.3 \times 10^{-15}$ on posterior means, $3.9 \times 10^{-15}$ on
variances, $1.2 \times 10^{-14}$ on scores, i.e.\ double-precision
rounding accumulated along the recursion. An independent implementation,
written from scratch against the recursion specification alone,
reproduces the same agreement (Appendix~\ref{app:repro}). One qualitative
detail from a worked $K = 3$ instance: at $x = (1.2, -0.4, 0.8)$ the score
component $S_1$ is positive although $x_1 < 0$, because the neighbours
pull the posterior mean of the middle frame positive. No per-frame
marginal score can produce this; it is inter-frame evidence sharing in
numbers.

\section{The bulk in closed form, and the price of locality}
\label{sec:g-bulk}\label{sec:g-locality}\label{sec:bp-truncation}

For a long chain the interior (bulk) is translation invariant and can be
solved exactly. By Proposition~\ref{prop:g-posterior} the posterior
precision in the bulk is tridiagonal Toeplitz with
\begin{equation}
  \boxed{\;
  J_d \;=\; \frac{e^{-2t}}{\Deltat} + \frac{1+\alpha^2}{\sigma_\eta^2},
  \qquad
  \beta \;=\; -\frac{\alpha}{\sigma_\eta^2},
  \qquad \sigma_\eta^2 = 1 - \alpha^2 .
  \;}
  \label{eq:g-bulk-params}
\end{equation}
The diagonal $J_d$ is the sum of the evidence precision $e^{-2t}/\Deltat$
(divergent as $1/2t$ at small $t$, vanishing at large $t$) and the
constant prior precision; the coupling $\beta$ is pure prior. The
dimensionless ratio $|\beta|/J_d \in (0, \tfrac{1}{2})$ controls
everything below.

\begin{theorem}[Exact bulk variance and geometric covariance]
\label{thm:g-bulk-cov}
In the bulk of a long chain, the posterior variance and covariance are
\begin{equation}
  \boxed{\;
  V \;=\; \frac{1}{\sqrt{J_d^2 - 4\beta^2}},
  \qquad
  \big( J^{-1} \big)_{i,\, i+d} \;=\; q^{\,d}\; V,
  \qquad
  q \;=\; \frac{J_d - \sqrt{J_d^2 - 4\beta^2}}{2\,|\beta|} \in (0, 1).
  \;}
  \label{eq:g-bulk-cov}
\end{equation}
\end{theorem}

\begin{toolbox}{Transfer-matrix inversion of a tridiagonal Toeplitz matrix}
Seek the inverse row in geometric form,
$(J^{-1})_{i, i+d} = C q^{\,d}$ for $d \geq 0$, symmetric in $\pm d$. The
defining relation $\sum_j J_{ij} (J^{-1})_{j, i+d} = \delta_{0d}$ reads,
for $d \geq 1$,
\begin{equation*}
  \beta\, C q^{\,d-1} + J_d\, C q^{\,d} + \beta\, C q^{\,d+1} = 0
  \quad\Longleftrightarrow\quad
  \beta q^2 + J_d\, q + \beta = 0,
\end{equation*}
a quadratic whose root inside the unit disc is
$q = (J_d - \sqrt{J_d^2 - 4\beta^2})/(2|\beta|)$ (for $\alpha > 0$ all
inverse entries are positive; for $\alpha < 0$ they alternate in sign with
the same modulus). The $d = 0$ relation fixes the prefactor:
$J_d C + 2\beta C q = 1$, and since
$2\beta q = -(J_d - \sqrt{J_d^2 - 4\beta^2})$,
$C = 1/\sqrt{J_d^2 - 4\beta^2} = V$. Existence of the real square root
requires $J_d > 2|\beta|$, which always holds here:
\begin{equation*}
  J_d - 2|\beta|
  = \frac{e^{-2t}}{\Deltat} + \frac{1 + \alpha^2 - 2|\alpha|}{\sigma_\eta^2}
  = \frac{e^{-2t}}{\Deltat} + \frac{(1 - |\alpha|)^2}{\sigma_\eta^2}
  \;>\; 0
\end{equation*}
for every $|\alpha| < 1$ and $t > 0$. This positivity margin is also what
guarantees the belief-propagation fixed point exists for every coupling
and diffusion time (Appendix~\ref{app:amp}). \hfill$\square$
\end{toolbox}

The number $q(\alpha, t)$ measures how far evidence travels along the
chain: the posterior correlation length is $\xi = 1/\log(1/q)$. Its limits
are
\begin{equation}
  \begin{aligned}
  t \to 0:&\quad q \approx \frac{|\beta|}{J_d} \to 0
  &&(\text{each frame pinned by its own observation});\\
  t \to \infty:&\quad q \to |\alpha|
  &&(\text{posterior} \to \text{prior}),
  \end{aligned}
  \label{eq:g-q-limits}
\end{equation}
the second because
$q \to \big((1+\alpha^2) - (1-\alpha^2)\big)/(2|\alpha|) = |\alpha|$ when
the evidence term vanishes. As diffusion time runs, the posterior
correlation length climbs monotonically from $0$ to the prior's own
$1/\log(1/|\alpha|)$.

\subsection{The locality law}

RQ4 asks how rapidly the influence of remote observations decays. Define
the local estimator of radius $r$: the posterior mean of $a_k$ given only
the observations within distance $r$,
$m_k^{(r)} = \E[a_k \,|\, x_{k-r}, \dots, x_{k+r}]$. Two facts make the
comparison clean. Any contiguous window of a stationary AR(1) chain is
itself a stationary AR(1) chain, so $m^{(r)}_k$ is exact on its window;
the only approximation is the truncation of range. And both estimators are
linear in $x$, so the error $m_k - m_k^{(r)} = e^\T x$ has a closed-form
RMS over the data law, $\mathrm{RMS}_r = \sqrt{e^\T \Sigt e}$, a
deterministic number with no sampling involved.

\begin{proposition}[Geometric locality: influence coefficients]
\label{thm:g-locality}
In the bulk, the influence of the observation at distance $d$ on the exact
posterior mean is $(J^{-1})_{k,k\pm d} \propto q(\alpha,t)^{\,d}$, exactly,
by Theorem~\ref{thm:g-bulk-cov}. Truncating the evidence at radius $r$
therefore discards a leading term of order $q^{\,r+1}$.
\end{proposition}

What Proposition~\ref{thm:g-locality} controls is the size of the
\emph{discarded} evidence, and that part is exact. It does not by itself
prove that the radius-$r$ estimator's error inherits the rate: $m^{(r)}_k$
is the posterior mean of a \emph{re-normalised} model on the window, so
truncation changes the boundary conditions as well as removing the leading
coefficient, and bounding the difference between those two effects is not
done here.

The rate is instead confirmed numerically, and closely. Fitted slopes of
$\log\mathrm{RMS}_r$ over $r = 2, \dots, 13$ at the centre of a $K = 121$
chain match $\log q$ within $0.2\%$ at $t \in \{0.1, 0.5, 2.0\}$
($\alpha = 0.8$); see Figure~\ref{fig:g-locality}, where the dashed lines
carry the predicted slope with no fitting. The claim this thesis makes is
therefore: an exact geometric law for the influence coefficients, and a
measured geometric law with the same rate for the windowed estimator's
error, over the stated range of $r$. A two-sided bound establishing the
second from the first is left open.

\begin{figure}[t]\centering
\includegraphics[width=\linewidth]{fig_local_vs_full.pdf}
\caption[Geometric decay of the windowed estimator's error]{Locality of the joint score ($\alpha = 0.8$, $K = 121$,
deterministic RMS over the data law). Left: exact truncation error of the
radius-$r$ local estimator (dots) at three diffusion times; the dashed
lines have the predicted slopes $q(\alpha,t)^r$ from
\eqref{eq:g-bulk-cov}, with no fitting. Right: the same error at fixed
radius as a function of $t$: it vanishes at both ends and peaks at
intermediate $t$, where the band of $\Qt$ is fullest.}
\label{fig:g-locality}
\end{figure}

\subsection{Truncating the matrix versus truncating the inference}

A truncated score can be built in two inequivalent ways.
\emph{Truncate the answer}: band the matrix,
$\widehat S = -Q_t^{(b)} x$ with entries beyond $|i-j| \leq b$ zeroed.
\emph{Truncate the algorithm}: let messages travel at most $r$ hops, i.e.\
run exact inference on the window $x_{k-r}, \dots, x_{k+r}$ (the local
estimator above). Both are linear in $x$, so their RMS errors are again
deterministic. Table~\ref{tab:g-truncation} compares them at equal locality
budget ($b = r = 1$, $K = 40$, $\alpha = 0.8$):
\begin{center}
\begin{minipage}{\linewidth}\centering
\captionof{table}[Truncating the answer against truncating the algorithm]{Truncating the answer against truncating the algorithm, at equal locality budget ($b = r = 1$, $K = 40$, $\alpha = 0.8$). Both estimators are linear in $x$, so both errors are deterministic rather than sampled. The cost of locality peaks at intermediate diffusion times and the ordering reverses nowhere.}
\label{tab:g-truncation}
\begin{tabular}{@{}lcccc@{}}
\toprule
relative RMS score error & $t = 0.05$ & $t = 0.3$ & $t = 1$ & $t = 3$ \\
\midrule
banded matrix ($b = 1$)      & $0.210$ & $0.302$ & $0.135$ & $0.0035$ \\
truncated inference ($r = 1$) & $0.123$ & $0.235$ & $0.130$ & $0.0035$ \\
\bottomrule
\end{tabular}
\end{minipage}
\end{center}
Two observations. The cost of locality is non-monotone in $t$ and peaks at
intermediate times, consistent with Section~\ref{sec:g-tridiag-loss}. And
at small $t$ truncating the algorithm beats truncating the answer
($12.3\%$ against $21.0\%$ at $t = 0.05$): the windowed inference solves a
correctly normalised problem on its window, while zeroing entries of an
inverse produces a matrix that is the inverse of nothing and misstates the
diagonal. The two coincide at large $t$, where the off-diagonal content
vanishes. Either error is measured to decay geometrically in the allowed
range at the rate $q(\alpha, t)$ of Proposition~\ref{thm:g-locality}.

\section{Status of the results}\label{sec:g-status}

Classified in the vocabulary of Section~\ref{sec:method}.
Exact, by derivation: the score formulas
\eqref{eq:g-score-matrix} and \eqref{eq:g-tweedie}, the spectral form
\eqref{eq:g-Qt-spec}, the band-fill law \eqref{eq:g-bandfill} as a
small-$t$ asymptotic, the large-$t$ law \eqref{eq:g-larget}, tree
exactness and Gaussian closure of BP, the Kalman identification, and the
bulk formulas \eqref{eq:g-bulk-params}--\eqref{eq:g-bulk-cov}, including
the $q^d$ decay of the influence coefficients. Numerical, within stated
grids and tolerances: the agreement of the two score routes, the audited
band-fill coefficients, the locality slopes of the \emph{windowed
estimator} (whose rate is measured rather than derived, see
Proposition~\ref{thm:g-locality}), and the truncation table. Interpretation, limited to
this model: within the Gaussian benchmark, a local score computation with
receptive field $r \gtrsim \xi \log(1/\varepsilon)$,
$\xi = 1/\log(1/q)$, achieves relative error $\varepsilon$, and the
requirement is most demanding at intermediate diffusion times; whether
trained score networks on real sequence data behave accordingly is not
tested here.
